\documentclass[11pt,a4paper]{article}
\usepackage[utf8]{inputenc}
\usepackage[spanish]{babel}
\usepackage{amsmath,amsfonts,amssymb,graphicx,booktabs,hyperref}
\title{\textbf{Proyecto Integrador: Modelado Predictivo para la Detección Temprana de la Deserción Universitaria}}
\author{Máquina Virtual de Minería de Datos}
\date{\today}
\begin{document}
\maketitle
\begin{abstract}
Este proyecto presenta la implementación de una solución analítica integral para predecir el riesgo de deserción en estudiantes de educación superior mediante minería de datos...
\end{abstract}
\section{Introducción}
El abandono en la educación superior representa un reto multidimensional...
\section{Definición del Problema}
Identificar tempranamente a estudiantes en riesgo de deserción antes del segundo ciclo académico...
\section{Base de Datos Seleccionada}
Se utilizó el benchmark público \textit{Predict Students Dropout and Academic Success} (UCI Machine Learning Repository)...
\section{Preparación y Calidad de los Datos}
Imputación de valores atípicos, validación de integridad y escalamiento Min-Max...
\section{Técnicas y Algoritmos Seleccionados}
Se contrastaron tres familias metodológicas: Árboles CART, Redes Neuronales MLP y Regresión Logística...
\section{Implementación en R}
El flujo se programó de forma reproducible con las librerías \texttt{rpart} y \texttt{nnet}...
\section{Resultados Experimentales}
El modelo de red neuronal alcanzó una precisión superior al 85\% en la muestra de prueba independiente...
\section{Evaluación de Modelos}
Matriz de confusión ponderada por el coste de no detectar a un estudiante en riesgo...
\section{Interpretación de Patrones}
Las cuotas al día y los créditos aprobados en el primer semestre mostraron el mayor peso sináptico discriminante...
\section{Conclusiones y Recomendaciones}
La minería de datos permite anticipar el 88\% de los casos de abandono antes de la matrícula del segundo período...
\section{Referencias Bibliográficas}
Romero, C., \& Ventura, S. (2010). Educational data mining: a review of the state of the art.
\end{document}
